\section{Diseño del Sistema}

En este cap\'itulo se presenta la arquitectura del sistema, las reglas de
gobernanza de datos, los flujos de negocio, las m\'aquinas de estado de las entidades
principales, el dise\~no de la API y el dise\~no de interfaces de usuario.

% -----------------------------------------------------------
\subsection{Arquitectura del Sistema}
% -----------------------------------------------------------

\subsubsection{Patr\'on Arquitect\'onico}

La arquitectura de la plataforma sigue el patr\'on de \textbf{Arquitectura en Capas}
(\textit{Layered Architecture}) combinado con el modelo de despliegue \textbf{serverless}
de Vercel. Las capas est\'an desacopladas mediante contratos de API bien definidos, lo que
facilita el mantenimiento independiente de cada componente y la escalabilidad horizontal
sin cambios en la l\'ogica de negocio. La Tabla~\ref{tab:capas_arq} describe las seis
capas de la arquitectura y sus tecnolog\'ias asociadas.

\begin{table}[ht]
\caption{Vista de Capas de la Arquitectura del Sistema}
\label{tab:capas_arq}
\centering
\small
\setlength{\tabcolsep}{3pt}
\begin{tabular}{>{\bfseries\raggedright\arraybackslash}p{3cm}
                >{\raggedright\arraybackslash}p{3.5cm}
                >{\raggedright\arraybackslash}p{9.5cm}}
\toprule
\thead{Capa} & \thead{Tecnología} & \thead{Descripción} \\
\midrule
Presentaci\'on &
  Next.js 14 + React 18 + TypeScript &
  SSR y CSR seg\'un la naturaleza de cada vista. Componentes React con estado local
  (\texttt{useState}/\texttt{useReducer}) y global (Zustand). Tailwind CSS para estilos.
  PWA accesible desde cualquier navegador moderno sin instalaci\'on nativa. \\
\rowcolor{altrow}
API / BFF &
  Next.js API Routes (Serverless en Vercel) &
  \textit{Backend For Frontend} implementado como funciones serverless. Valida JWT de
  Supabase Auth, aplica reglas de negocio y delega operaciones a Prisma ORM.
  Timeout m\'aximo: 60~s (plan Pro). Edge Network de Vercel ($>$100 ubicaciones). \\
Dominio / ORM &
  Prisma ORM 5 + Prisma Accelerate &
  El schema de Prisma define todas las entidades y relaciones. Genera un cliente
  TypeScript tipado que previene SQL Injection por dise\~no. Prisma Accelerate
  act\'ua como proxy de connection pooling entre funciones serverless y PostgreSQL \cite{mdn_pwa}. \\
\rowcolor{altrow}
Persistencia &
  Supabase (PostgreSQL 15 + Auth + Realtime) &
  Base de datos relacional con RLS en todas las tablas sensibles. Supabase Auth
  gestiona sesiones con JWT RS256. Supabase Realtime (LISTEN/NOTIFY) distribuye
  eventos a clientes suscritos sin polling \cite{supabase_license}. \\
Servicios Persistentes &
  VPS (Ubuntu 22.04 LTS, Node.js 20 LTS) &
  Servicios de ejecuci\'on continua: procesamiento asíncrono, jobs de expiraci\'on de
  mesas y servicios WebSocket de respaldo. Conexi\'on directa a PostgreSQL (sin pooler,
  conexiones de larga duraci\'on). \\
\rowcolor{altrow}
Servicios Externos &
  Email (Resend/Sendgrid) &
  Env\'io de comprobantes digitales y notificaciones. La capa de pago con tarjeta
  queda fuera del alcance de la v1.0 (ver LM-13). \\
\bottomrule
\end{tabular}
\end{table}

\subsubsection{Patrones de Diseño Aplicados}

La Tabla~\ref{tab:patrones} resume los patrones de dise\~no aplicados en la
arquitectura del sistema y el beneficio concreto que aporta cada uno.

\begin{table}[ht]
\caption{Patrones de Diseño Aplicados en el Sistema}
\label{tab:patrones}
\centering
\small
\setlength{\tabcolsep}{3pt}
\begin{tabular}{>{\bfseries\raggedright\arraybackslash}p{3.5cm}
                >{\raggedright\arraybackslash}p{6.0cm}
                >{\raggedright\arraybackslash}p{6.5cm}}
\toprule
\thead{Patrón} & \thead{Aplicación en el sistema} & \thead{Beneficio} \\
\midrule
Repository Pattern &
  Prisma ORM abstrae el acceso a PostgreSQL. Las API Routes nunca escriben SQL directo. &
  Desacoplamiento entre l\'ogica de negocio y persistencia.
  Facilita pruebas unitarias con mocks. \\
\rowcolor{altrow}
Backend For Frontend &
  Las API Routes de Next.js act\'uan como BFF: una API dise\~nada espec\'ificamente
  para el frontend de la plataforma. &
  Centraliza validaciones de autenticaci\'on y autorizaci\'on. Evita exponer la base
  de datos directamente al cliente. \\
Optimistic UI &
  El cliente actualiza el estado local al agregar una orden sin esperar confirmaci\'on
  del servidor. Rollback si el servidor rechaza. &
  Percepci\'on de respuesta instant\'anea para el comensal. \\
\rowcolor{altrow}
Observer (Realtime) &
  Supabase Realtime: los clientes se suscriben a canales (mesas, \'ordenes) y reciben
  notificaciones de cambios sin polling. &
  Actualizaciones en tiempo real para el Chef y los comensales sin impacto
  en el rendimiento del servidor \cite{nextjs}. \\
Finite State Machine &
  Las mesas virtuales y las cuentas tienen estados bien definidos con transiciones
  expl\'icitas controladas por el backend (ver Secci\'on~5.3). &
  Previene estados inconsistentes. Facilita el razonamiento sobre el flujo
  de liquidaci\'on. \\
\rowcolor{altrow}
Idempotency Key &
  Cada operaci\'on de pago genera una clave \'unica de idempotencia almacenada
  antes del procesamiento. &
  Previene cobros duplicados si la red falla y el cliente reintenta (RT-01, RT-02). \\
Row Level Security &
  Las pol\'iticas RLS de PostgreSQL garantizan acceso solo a los datos del nivel
  organizacional del usuario. &
  Defensa en profundidad: aunque la API tenga un bug de autorizaci\'on, la base de
  datos rechaza el acceso no autorizado. \\
\bottomrule
\end{tabular}
\end{table}

\subsubsection{Diagrama de Despliegue}

La Figura~\ref{fig:deployment} muestra la topolog\'ia de despliegue del sistema, con los nodos f\'isicos
y virtuales, los artefactos desplegados en cada uno y los protocolos de comunicación.

\begin{figure}[ht]
\centering
\begin{usecasebox}{Diagrama de Despliegue de la Plataforma}
\centering
\includegraphics[width=0.95\textwidth]{architecture/rendered/deployment.pdf}
\end{usecasebox}
\caption{Diagrama de despliegue de la plataforma propuesta (topología de nodos y protocolos).}
\label{fig:deployment}
\end{figure}

\subsubsection{Diagrama de Clases UML --- Entidades Principales}

El diagrama de clases de la Figura~\ref{fig:clases_uml} modela las entidades principales del sistema y
sus relaciones, siguiendo la notaci\'on UML~2.0 [documento de apoyo]. Las
asociaciones de composici\'on (rombo relleno) indican que el objeto hijo no puede
existir sin el padre; las asociaciones simples (l\'inea) indican navegabilidad
entre entidades independientes.

\begin{figure}[ht]
\centering
\begin{usecasebox}{Diagrama de Clases UML --- Entidades Principales}
\centering
\includegraphics[width=0.95\textwidth]{architecture/rendered/classes.pdf}
\end{usecasebox}
\caption{Diagrama de clases UML de las entidades principales del sistema.
Las composiciones ($\blacklozenge$) indican dependencia existencial; las agregaciones
($\lozenge$) y asociaciones (---) indican relaciones navegables independientes.}
\label{fig:clases_uml}
\end{figure}

\subsubsection{Diagrama de Secuencia --- Flujo de Activación de Mesa Virtual}

La Figura~\ref{fig:seq_activacion} muestra el diagrama de secuencia del flujo de activaci\'on de mesa virtual
por parte del Anfitr\'ion, que es el flujo m\'as cr\'itico del sistema por involucrar
la generaci\'on del JWT, la validaci\'on del QR y la creaci\'on de la sesi\'on.

\begin{figure}[ht]
\centering
\begin{usecasebox}{Diagrama de Secuencia --- Activación de Mesa Virtual (Flujo Anfitrión)}
\centering
\includegraphics[width=0.95\textwidth]{casos_uso/rendered/CU_01.pdf}
\end{usecasebox}
\caption{Diagrama de secuencia del flujo de activaci\'on de mesa virtual por el Anfitr\'ion.
Los pasos 11 son at\'omicos (transacci\'on PostgreSQL) para prevenir activaciones concurrentes (RO-07).}
\label{fig:seq_activacion}
\end{figure}

\FloatBarrier

\subsubsection{CU-04 y CU-05}
\noindent (Movidos a la Secci\'on 4.4: Casos de Uso del Sistema).

\FloatBarrier

\FloatBarrier

% -----------------------------------------------------------
\subsection{Diseño de Base de Datos}
% -----------------------------------------------------------

\subsubsection{Modelo Entidad-Relación (MER)}

La base de datos relacional del sistema se implementa sobre PostgreSQL~15
a trav\'es de Supabase, gestionada mediante el esquema de Prisma~ORM~5. El dise\~no
sigue el modelo relacional normalizado con integridad referencial garantizada mediante
claves for\'aneas, transacciones ACID y pol\'iticas de Row Level Security (RLS) en
todas las tablas sensibles \cite{supabase_license, mdn_pwa}.

El esquema se organiza en cuatro dominios funcionales claramente delimitados:
(1)~\textbf{Dominio organizacional} (Chain, Zone, Restaurant), que modela la
jerarqu\'ia de tres niveles; (2)~\textbf{Dominio de personal} (ChainStaff, Staff),
que gestiona los roles de empleados; (3)~\textbf{Dominio operativo} (Table, Session,
Order, OrderItem, Payment, PaymentAllocation), que soporta el ciclo completo de
servicio; y (4)~\textbf{Dominio de cat\'alogo} (MenuTemplate, TemplateCategory,
TemplateItem, Category, MenuItem, RestaurantMenuOverride, ZoneMenuOverride), que
implementa el sistema de men\'us con soporte de plantillas compartidas y overrides
por zona o sucursal.

La Figura~\ref{fig:mer} presenta el Modelo Entidad-Relaci\'on simplificado con las 19~tablas del
esquema y sus relaciones principales.

\begin{figure}[ht]
\begin{usecasebox}{Modelo Entidad-Relación --- Base de Datos (PostgreSQL 15)}
\centering
\includegraphics[width=0.95\textwidth]{architecture/rendered/mer_global.pdf}
\end{usecasebox}
\caption{Modelo Entidad-Relaci\'on del sistema. Las composiciones ($\blacklozenge$)
indican dependencia existencial. El dominio de cat\'alogo implementa herencia por
override: TemplateItem define el precio base; ZoneMenuOverride y
RestaurantMenuOverride permiten personalizar precio, disponibilidad y visibilidad
por nivel organizacional.}
\label{fig:mer}
\end{figure}

\subsubsection{Diccionario de Datos --- Tablas Principales}

Las tablas siguientes documentan las entidades m\'as relevantes del esquema. El
diccionario completo se encuentra en el repositorio del proyecto bajo
\texttt{prisma/schema.prisma}.

La Tabla~\ref{tab:dd_chain} describe la entidad ra\'iz de la jerarqu\'ia organizacional.

\begin{table}[ht]
\caption{Diccionario de Datos --- Tabla \texttt{Chain} (Cadena Restaurantera)}
\label{tab:dd_chain}
\centering
\small
\setlength{\tabcolsep}{3pt}
\begin{tabular}{>{\ttfamily\raggedright\arraybackslash}p{3cm}
                >{\raggedright\arraybackslash}p{2.5cm}
                >{\centering\arraybackslash}p{1.5cm}
                >{\centering\arraybackslash}p{1.5cm}
                >{\raggedright\arraybackslash}p{7.5cm}}
\toprule
\thead{Campo} & \thead{Tipo} & \thead{PK/FK} & \thead{Nulo} & \thead{Descripción / Restricción} \\
\midrule
id & text & PK & No & Identificador \'unico generado por CUID (Prisma). \\
\rowcolor{altrow}
name & text & --- & No & Nombre de la cadena restaurantera. \\
logoUrl & text & --- & S\'i & URL p\'ublica del logo almacenado en Supabase Storage. \\
\rowcolor{altrow}
currency & text & --- & No & Moneda por defecto. DEFAULT: \texttt{'MXN'}. \\
taxRate & float & --- & No & Tasa de IVA aplicable. DEFAULT: \texttt{16.0} (\%). \\
\rowcolor{altrow}
suggestedTips & text & --- & No & Propinas sugeridas separadas por coma. DEFAULT: \texttt{'10,15,20'}. \\
createdAt & timestamp & --- & No & Timestamp de creaci\'on (UTC). DEFAULT: \texttt{NOW()}. \\
\rowcolor{altrow}
updatedAt & timestamp & --- & No & Timestamp de \'ultima modificaci\'on. Actualizado autom\'aticamente por Prisma. \\
\bottomrule
\end{tabular}
\end{table}

La Tabla~\ref{tab:dd_table} describe la entidad mesa f\'isica, central en el flujo operativo del sistema.

\begin{table}[ht]
\caption{Diccionario de Datos --- Tabla \texttt{Table} (Mesa Física)}
\label{tab:dd_table}
\centering
\small
\setlength{\tabcolsep}{3pt}
\begin{tabular}{>{\ttfamily\raggedright\arraybackslash}p{3cm}
                >{\raggedright\arraybackslash}p{2.5cm}
                >{\centering\arraybackslash}p{1.5cm}
                >{\centering\arraybackslash}p{1.5cm}
                >{\raggedright\arraybackslash}p{7.5cm}}
\toprule
\thead{Campo} & \thead{Tipo} & \thead{PK/FK} & \thead{Nulo} & \thead{Descripción / Restricción} \\
\midrule
id & text & PK & No & Identificador \'unico (CUID). \\
\rowcolor{altrow}
restaurantId & text & FK & No & Referencia a \texttt{Restaurant.id}. CASCADE on delete. \\
number & integer & --- & No & N\'umero interno de mesa asignado por el restaurante. \\
\rowcolor{altrow}
capacity & integer & --- & No & Capacidad m\'axima de comensales. \\
qrCode & text & --- & No & Token JWT de un solo uso generado al inicializar la mesa. Se invalida al ser escaneado. \\
\rowcolor{altrow}
status & TableStatus & --- & No & Enum: \texttt{DISPONIBLE | OCUPADA | CERRANDO | SUCIA}. DEFAULT: \texttt{DISPONIBLE}. \\
posX & float & --- & No & Posici\'on horizontal en el plano visual (dashboard). DEFAULT: \texttt{80}. \\
\rowcolor{altrow}
posY & float & --- & No & Posici\'on vertical en el plano visual. DEFAULT: \texttt{80}. \\
shape & text & --- & No & Forma visual: \texttt{'rect'} o \texttt{'circle'}. DEFAULT: \texttt{'rect'}. \\
\rowcolor{altrow}
joinCode & text & --- & S\'i & C\'odigo p\'ublico de uni\'on generado al activar la mesa. Rotado al cerrar sesi\'on. \\
parentTableId & text & FK & S\'i & Auto-referencia para mesa combinada (tabla unida). NULL si es mesa independiente. \\
\bottomrule
\end{tabular}
\end{table}

La Tabla~\ref{tab:dd_session} describe la sesi\'on de comensal, entidad que vincula un usuario con una mesa activa.

\begin{table}[ht]
\caption{Diccionario de Datos --- Tabla \texttt{Session} (Sesión de Comensal)}
\label{tab:dd_session}
\centering
\small
\setlength{\tabcolsep}{3pt}
\begin{tabular}{>{\ttfamily\raggedright\arraybackslash}p{3cm}
                >{\raggedright\arraybackslash}p{2.5cm}
                >{\centering\arraybackslash}p{1.5cm}
                >{\centering\arraybackslash}p{1.5cm}
                >{\raggedright\arraybackslash}p{7.5cm}}
\toprule
\thead{Campo} & \thead{Tipo} & \thead{PK/FK} & \thead{Nulo} & \thead{Descripción / Restricción} \\
\midrule
id & text & PK & No & Identificador \'unico (CUID). \\
\rowcolor{altrow}
tableId & text & FK & No & Referencia a \texttt{Table.id}. \\
guestName & text & --- & No & Nombre del comensal ingresado al unirse a la mesa. \\
\rowcolor{altrow}
pax & integer & --- & No & N\'umero de personas en la sesi\'on. \\
isActive & boolean & --- & No & Indica si la sesi\'on est\'a activa. DEFAULT: \texttt{true}. \\
\rowcolor{altrow}
isHost & boolean & --- & No & Indica si el comensal es el anfitr\'ion de la mesa. DEFAULT: \texttt{false}. Solo un comensal por mesa puede tener \texttt{isHost = true}. \\
createdAt & timestamp & --- & No & Timestamp de ingreso del comensal a la mesa. \\
\rowcolor{altrow}
closedAt & timestamp & --- & S\'i & Timestamp de cierre de sesi\'on. NULL mientras la mesa est\'a activa. \\
\bottomrule
\end{tabular}
\end{table}

La Tabla~\ref{tab:dd_order} describe la entidad orden, central en el flujo de pedidos.

\begin{table}[ht]
\caption{Diccionario de Datos --- Tabla \texttt{Order} (Orden)}
\label{tab:dd_order}
\centering
\small
\setlength{\tabcolsep}{3pt}
\begin{tabular}{>{\ttfamily\raggedright\arraybackslash}p{3cm}
                >{\raggedright\arraybackslash}p{2.5cm}
                >{\centering\arraybackslash}p{1.5cm}
                >{\centering\arraybackslash}p{1.5cm}
                >{\raggedright\arraybackslash}p{7.5cm}}
\toprule
\thead{Campo} & \thead{Tipo} & \thead{PK/FK} & \thead{Nulo} & \thead{Descripción / Restricción} \\
\midrule
id & text & PK & No & Identificador \'unico (CUID). \\
\rowcolor{altrow}
restaurantId & text & FK & No & Referencia a \texttt{Restaurant.id}. \\
tableId & text & FK & No & Referencia a \texttt{Table.id}. \\
\rowcolor{altrow}
status & OrderStatus & --- & No & Enum: \texttt{PENDING | IN\_PROGRESS | READY | DELIVERED | CANCELLED}. DEFAULT: \texttt{PENDING}. \\
guestName & text & --- & S\'i & Nombre del comensal que gener\'o la orden. Desnormalizado para trazabilidad hist\'orica. \\
\rowcolor{altrow}
deliveredAt & timestamp & --- & S\'i & Timestamp de entrega al comensal. NULL hasta que el comensal confirma recepci\'on. \\
\bottomrule
\end{tabular}
\end{table}

La Tabla~\ref{tab:dd_payment} describe la entidad de pago con sus campos de control de concurrencia.

\begin{table}[ht]
\caption{Diccionario de Datos --- Tabla \texttt{Payment} (Pago)}
\label{tab:dd_payment}
\centering
\small
\setlength{\tabcolsep}{3pt}
\begin{tabular}{>{\ttfamily\raggedright\arraybackslash}p{3cm}
                >{\raggedright\arraybackslash}p{2.5cm}
                >{\centering\arraybackslash}p{1.5cm}
                >{\centering\arraybackslash}p{1.5cm}
                >{\raggedright\arraybackslash}p{7.5cm}}
\toprule
\thead{Campo} & \thead{Tipo} & \thead{PK/FK} & \thead{Nulo} & \thead{Descripción / Restricción} \\
\midrule
id & text & PK & No & Identificador \'unico (CUID). \\
\rowcolor{altrow}
restaurantId & text & FK & No & Referencia a \texttt{Restaurant.id}. \\
tableId & text & FK & No & Referencia a \texttt{Table.id}. \\
\rowcolor{altrow}
sessionId & text & FK & No & Referencia a \texttt{Session.id}. Identifica al comensal pagador. \\
orderId & text & FK & S\'i & Referencia opcional a \texttt{Order.id}. NULL si el pago cubre m\'ultiples \'ordenes. \\
\rowcolor{altrow}
status & PaymentStatus & --- & No & Enum: \texttt{PENDING | COMPLETED | FAILED | REFUNDED}. DEFAULT: \texttt{PENDING}. \\
method & PaymentMethod & --- & No & Enum: \texttt{CARD | CASH | TRANSFER}. DEFAULT: \texttt{CARD}. \\
\rowcolor{altrow}
splitMode & PaymentSplitMode & --- & No & Enum: \texttt{FULL | BY\_ITEM | EQUAL | CUSTOM}. DEFAULT: \texttt{FULL}. \\
splitCount & integer & --- & No & N\'umero total de partes en la divisi\'on. CHECK $\geq 1$. \\
\rowcolor{altrow}
paxPaid & integer & --- & No & N\'umero de personas cubierto por este pago. CHECK $\geq 1$. \\
subtotal & float & --- & No & Monto antes de propina e impuestos. CHECK $\geq 0$. \\
\rowcolor{altrow}
tipRate & float & --- & No & Porcentaje de propina seleccionado. DEFAULT: \texttt{0}. \\
tipAmount & float & --- & No & Monto calculado de propina. CHECK $\geq 0$. \\
\rowcolor{altrow}
totalAmount & float & --- & No & Subtotal + propina. CHECK $\geq 0$. \\
amountPaid & float & --- & No & Monto efectivamente cobrado (puede diferir en pagos parciales). CHECK $\geq 0$. \\
\rowcolor{altrow}
paidAt & timestamp & --- & S\'i & Timestamp de confirmaci\'on del pago. NULL hasta completarse. \\
\bottomrule
\end{tabular}
\end{table}

La Tabla~\ref{tab:dd_menutemplate} describe el sistema de plantillas de men\'u compartido entre sucursales.

\begin{table}[ht]
\caption{Diccionario de Datos --- Sistema de Plantillas de Menú (MenuTemplate, TemplateCategory, TemplateItem)}
\label{tab:dd_menutemplate}
\centering
\small
\setlength{\tabcolsep}{3pt}
\begin{tabular}{>{\raggedright\arraybackslash}p{3cm}
                >{\ttfamily\raggedright\arraybackslash}p{3cm}
                >{\raggedright\arraybackslash}p{2.5cm}
                >{\raggedright\arraybackslash}p{7.5cm}}
\toprule
\thead{Tabla} & \thead{Campo clave} & \thead{Tipo} & \thead{Descripción} \\
\midrule
MenuTemplate & id, chainId (FK) & text & Plantilla maestra de men\'u definida a nivel cadena. \texttt{isDefault=true} indica la plantilla activa. \\
\rowcolor{altrow}
TemplateCategory & id, templateId (FK) & text & Categor\'ia de la plantilla (ej. Entradas, Bebidas). Campo \texttt{order} controla el orden de presentaci\'on. \\
TemplateItem & id, categoryId (FK) & text & \'Item base con \texttt{basePrice}. La estaci\'on de preparaci\'on se define con el enum \texttt{Station} (\texttt{COCINA}, \texttt{BARRA}). \\
\rowcolor{altrow}
ZoneMenuOverride & id, zoneId (FK), templateItemId (FK) & text & Override por zona: permite sobreescribir \texttt{customPrice}, \texttt{isSoldOut} e \texttt{isHidden} sin modificar la plantilla maestra. \\
RestaurantMenuOverride & id, restaurantId (FK), templateItemId (FK) & text & Override por sucursal: m\'axima granularidad de personalizaci\'on del men\'u. Hereda de ZoneMenuOverride si no tiene override propio. \\
\bottomrule
\end{tabular}
\end{table}

\subsubsection{Tipos Enumerados (ENUMs) del Esquema}

El esquema define los siguientes tipos enumerados a nivel de base de datos en
PostgreSQL, garantizando consistencia en los valores de estado de cada entidad.
La Tabla~\ref{tab:enums} los documenta con sus valores posibles.

\begin{table}[ht]
\caption{Tipos Enumerados (ENUMs) Definidos en el Esquema PostgreSQL}
\label{tab:enums}
\centering
\small
\setlength{\tabcolsep}{3pt}
\begin{tabular}{>{\ttfamily\bfseries\raggedright\arraybackslash}p{3.5cm}
                >{\raggedright\arraybackslash}p{6.0cm}
                >{\raggedright\arraybackslash}p{6.5cm}}
\toprule
\thead{Enum} & \thead{Valores} & \thead{Contexto de uso} \\
\midrule
TableStatus & DISPONIBLE, OCUPADA, CERRANDO, SUCIA & Estado de la mesa f\'isica. Controla el flujo de inicializaci\'on de sesi\'on y liberaci\'on. \\
\rowcolor{altrow}
OrderStatus & PENDING, IN\_PROGRESS, READY, DELIVERED, CANCELLED & Estado de la orden en el flujo del Chef. Transiciones controladas por el backend. \\
PaymentStatus & PENDING, COMPLETED, FAILED, REFUNDED & Estado del pago. PENDING es el \'unico estado mutable v\'ia SELECT FOR UPDATE. \\
\rowcolor{altrow}
PaymentMethod & CARD, CASH, TRANSFER & M\'etodo de liquidaci\'on. CARD es el valor por defecto. \\
PaymentSplitMode & FULL, BY\_ITEM, EQUAL, CUSTOM & Modalidad de divisi\'on de cuenta: total, por \'item, partes iguales o monto personalizado. \\
\rowcolor{altrow}
Station & COCINA, BARRA & Estaci\'on de preparaci\'on del \'item. Determina el destino de la orden en el KDS. \\
StaffRole & (Roles de empleados: ADMIN, GERENTE, MESERO, CHEF) & Rol del colaborador. Utilizado por las pol\'iticas RLS de Supabase para control de acceso. \\
\bottomrule
\end{tabular}
\end{table}

\subsubsection{Decisiones de Diseño Relevantes}

\noindent\textbf{Desnormalizaci\'on controlada.}
El campo \texttt{guestName} en \texttt{Order} y \texttt{priceAtTime} en
\texttt{OrderItem} son ejemplos de desnormalizaci\'on deliberada: almacenan el
nombre del comensal y el precio del producto \textit{en el momento de la orden},
preservando la trazabilidad hist\'orica incluso si el comensal cierra sesi\'on o
el gerente modifica el precio del men\'u con posterioridad.

\noindent\textbf{Sistema de overrides de men\'u en cascada.}
El dise\~no de cat\'alogo implementa un patr\'on de herencia por override en tres
niveles: la \texttt{MenuTemplate} define el men\'u base a nivel cadena;
\texttt{ZoneMenuOverride} permite personalizar precio y disponibilidad por zona;
y \texttt{RestaurantMenuOverride} a\~nade la mayor granularidad a nivel sucursal.
Esto permite que una cadena centralice su men\'u pero cada restaurante ajuste
precios localmente sin duplicar registros.

\noindent\textbf{Mesas combinadas (\textit{table joining}).}
El campo \texttt{parentTableId} en \texttt{Table} implementa una auto-referencia
que permite combinar m\'ultiples mesas f\'isicas en una sola sesi\'on operativa,
sin requerir una tabla de uni\'on separada. Una mesa con \texttt{parentTableId = null}
es independiente; las mesas con valor en ese campo son subcomponentes de la mesa padre.

\noindent\textbf{Control de concurrencia en pagos.}
La combinaci\'on de los campos \texttt{status}, \texttt{amountPaid} y
\texttt{totalAmount} en \texttt{Payment}, junto con el uso de
\texttt{SELECT FOR UPDATE} a nivel de transacci\'on PostgreSQL, implementa el
mecanismo anti-doble-cobro descrito en los riesgos RT-01 y RT-02.

\noindent\textbf{Rotaci\'on de c\'odigo QR.}
El campo \texttt{joinCode} en \texttt{Table} es generado y rotado en cada cierre
de sesi\'on por la acci\'on \texttt{closeTable} (Server Action en
\texttt{waiter.ts}), garantizando que un c\'odigo de uni\'on previo no pueda
ser reutilizado por comensales no autorizados.

\FloatBarrier

% -----------------------------------------------------------
\subsection{Gobernanza de Datos}
% -----------------------------------------------------------

\subsubsection{Estructura Organizacional}

La plataforma soporta tres niveles jerárquicos de organización para que las cadenas
restauranteras puedan estructurarse internamente:

\begin{itemize}[leftmargin=1.5em]
  \item Cadena Restaurantera
  \item Zona
  \item Sucursal
\end{itemize}

\subsubsection{Roles de Usuarios}

Los usuarios se dividen en dos grupos principales: colaboradores del restaurante
(pertenecientes a la operación) y consumidores (quienes consumen los productos y
servicios). Los roles de \textbf{Empleados} son: Administrador (Admin),
Gerente Zona, Gerente Sucursal, Mesero y Chef. Los roles de \textbf{Consumidores}
son: Anfitrión y Comensal.

Para garantizar la seguridad del sistema, cada grupo de comensales deberá contar
con un Anfitrión elegido entre ellos, quien activará la mesa virtual escaneando
el código QR que el mesero proporcione y creará una contraseña de acceso compartible
con los demás comensales.

Todo usuario colaborador debe estar ligado a alguno de los niveles de la estructura
organizacional y, dependiendo de su rol, tendrá acceso únicamente a su nivel y a
los niveles que este englobe.

\subsubsection{Reglas de Asignación de Roles a Niveles}

\begin{itemize}[leftmargin=1.5em]
  \item El usuario Admin únicamente existe a nivel Cadena de Restaurante.
  \item El usuario Gerente puede asignarse a nivel Zona (Gerente Zona) o a nivel
        Sucursal (Gerente Sucursal).
  \item Los usuarios Mesero y Chef únicamente pueden estar asignados a nivel Sucursal.
\end{itemize}

La Tabla~\ref{tab:permisos} presenta la matriz de permisos por rol del sistema.

\setstretch{1.2}
\setlength{\tabcolsep}{4pt}
\begin{longtable}{>{\raggedright\arraybackslash}p{10cm}
                >{\centering\arraybackslash}p{0.8cm}
                >{\centering\arraybackslash}p{1.1cm}
                >{\centering\arraybackslash}p{1.1cm}
                >{\centering\arraybackslash}p{1cm}
                >{\centering\arraybackslash}p{1.1cm}
                >{\centering\arraybackslash}p{0.8cm}}
\caption{Matriz de Permisos por Rol} \label{tab:permisos} \\
\toprule
\thead{Acción / Permiso} &
\thead{Chef} &
\thead{Comensal} &
\thead{Anfi-trión} &
\thead{Mesero} &
\thead{Gerente} &
\thead{Admin} \\
\midrule
\endfirsthead
\caption[]{TABLA \thetable{} (Continuación)} \\
\toprule
\thead{Acción / Permiso} &
\thead{Chef} &
\thead{Comensal} &
\thead{Anfi-trión} &
\thead{Mesero} &
\thead{Gerente} &
\thead{Admin} \\
\midrule
\endhead
\bottomrule
\endfoot
Dar de alta Zonas                          &   &   &   &   &   & X \\
\rowcolor{altrow}
Dar de alta usuarios con rol Admin         &   &   &   &   &   & X \\
Dar de alta usuarios con otros roles       &   &   &   &   & X & X \\
\rowcolor{altrow}
Acceso al Dashboard de Analytics           &   &   &   &   & X & X \\
Dar de alta Sucursales                     &   &   &   &   & X & X \\
\rowcolor{altrow}
Agregar, eliminar o editar menús/productos &   &   &   &   & X &   \\
Crear, editar y liberar mesas físicas      &   &   &   &   & X & X \\
\rowcolor{altrow}
Visualizar mesas físicas y características &   &   &   &   & X & X \\
Inicializar mesa virtual y su identificador&   &   &   & X & X & X \\
\rowcolor{altrow}
Activar mesa virtual y crear contraseña    &   &   & X &   &   &   \\
Visualizar y compartir contraseña de acceso&   & X & X &   &   &   \\
\rowcolor{altrow}
Generar órdenes individuales               &   & X & X &   &   &   \\
Generar órdenes compartidas                &   & X & X &   &   &   \\
\rowcolor{altrow}
Visualizar órdenes y su estatus            & X &   &   & X & X & X \\
Visualizar montos de cuentas por comensal  &   & X & X & X & X & X \\
\rowcolor{altrow}
Liquidar cuentas o aportar a mesa compartida&  & X & X &   &   &   \\
Visualizar estatus de liquidación          &   & X & X & X & X & X \\
\end{longtable}
\setstretch{1.5}

\noindent\textit{Nota: El Chef únicamente puede ver órdenes hasta el estado
EN\_REPARTO. El Mesero únicamente puede ver el cambio de EN\_REPARTO a ENTREGADA.}

\medskip

La Tabla~\ref{tab:alta_usuarios} define qué roles puede dar de alta cada tipo
de usuario del sistema.

\begin{table}[ht]
\caption{Reglas de Alta de Usuarios por Rol (RB-U)}
\label{tab:alta_usuarios}
\centering
\small
\setlength{\tabcolsep}{5pt}
\begin{tabular}{>{\raggedright\arraybackslash}p{6cm}
                >{\centering\arraybackslash}p{2cm}
                >{\centering\arraybackslash}p{2cm}
                >{\centering\arraybackslash}p{2cm}
                >{\centering\arraybackslash}p{2cm}
                >{\centering\arraybackslash}p{2cm}}
\toprule
\thead{¿Puede dar de alta?} &
\thead{Admin} &
\thead{Gerente Zona} &
\thead{Gerente Suc.} &
\thead{Mesero} &
\thead{Chef} \\
\midrule
Admin                & X & X & X & X & X \\
\rowcolor{altrow}
Gerente Zona         &   &   & X & X & X \\
Gerente Sucursal     &   &   &   & X & X \\
\rowcolor{altrow}
Mesero               &   &   &   &   &   \\
Chef                 &   &   &   &   &   \\
\bottomrule
\end{tabular}
\end{table}

% -----------------------------------------------------------
\subsection{Alcance del Sistema}
% -----------------------------------------------------------

\subsubsection{Dentro del Alcance}

\begin{itemize}[leftmargin=1.5em]
  \item Alta y gestión de mesas físicas para 1 sola área (1 piso) por sucursal.
  \item Gestión de mesas virtuales con flujo completo de activación, órdenes y liquidación.
  \item Módulo de órdenes individuales y mesa compartida para comensales.
  \item Tablero Kanban en tiempo real para el Chef.
  \item Dashboard de analytics para Gerentes y Admin.
  \item Gestión de menú (categorías, productos y variantes).
  \item Gestión de usuarios colaboradores con jerarquía de permisos.
  \item M\'odulo de liquidaci\'on con control de concurrencia y comprobante digital
        (sin agregador de pago v\'ia tarjeta en v1.0; ver LM-13).
\end{itemize}

\subsubsection{Fuera del Alcance}

\begin{itemize}[leftmargin=1.5em]
  \item Seccionamiento de mesas físicas por áreas o pisos múltiples.
  \item Módulo de cambio o recuperación de contraseña de mesa virtual en caso de olvido.
  \item Flujo de recuperación de contraseña para usuarios colaboradores.
  \item Base de datos analítica y clustering de datos para reducción de latencia en el Dashboard.
  \item Descarga de reportes XLSX por período de tiempo sobre transacciones.
  \item Visualización detallada para el mesero de qué usuarios faltan por liquidar
        y qué productos solicitó cada uno.
\end{itemize}

\subsubsection{Riesgos Identificados}

\begin{itemize}[leftmargin=1.5em]
  \item Dos comensales intentan pagar a un mismo comensal simultáneamente.
  \item Múltiples pagos a la mesa compartida se generan al mismo tiempo,
        superando el monto a liquidar.
  \item El anfitrión pierde acceso a la mesa virtual y no recuerda la contraseña,
        sin que ningún comensal haya ingresado previamente.
  \item Se comparten credenciales de la mesa virtual con personas no autorizadas.
  \item Dos comensales intentan activar la misma mesa virtual al mismo tiempo.
\end{itemize}

\subsubsection{Trabajo Futuro}

\begin{itemize}[leftmargin=1.5em]
  \item Modelo de machine learning para recomendación de platillos basado en
        historial del comensal.
  \item Flujo de cambio de contraseña de mesa virtual por parte del anfitrión.
  \item Botón para llamar al mesero desde la app del comensal.
  \item Notificaciones push al mesero para acciones requeridas por el usuario.
\end{itemize}

% -----------------------------------------------------------
\subsection{Flujos de Negocio}
% -----------------------------------------------------------

\subsubsection{Login de Usuario Colaborador}

Roles participantes: Admin, Gerente, Mesero, Chef. El usuario ingresa su correo
y contraseña en la pantalla de login. En el primer acceso el sistema detecta la
contraseña temporal y redirige al flujo de cambio de contraseña. Tras la
autenticación exitosa, el sistema redirige según el rol: Admin/Gerente al Dashboard
de Analytics; Mesero a la pantalla de mesas físicas de su sucursal; Chef al
tablero Kanban de órdenes de su sucursal.

\subsubsection{Alta de Zonas}

Roles participantes: Admin. El usuario accede al módulo de Zonas desde la barra
de navegación, puede visualizar zonas existentes y crear nuevas. Los campos
obligatorios son: nombre. Campo opcional: descripción.

\subsubsection{Alta de Sucursales}

Roles participantes: Admin, Gerente Zona. El Gerente Zona únicamente puede dar
de alta sucursales dentro de las zonas que tiene asignadas.

\subsubsection{Alta de Categorías y Productos}

Roles participantes: Admin, Gerente. Las categorías agrupan los productos del
menú. La creación está filtrada por las zonas o sucursales asignadas al usuario.
Un producto debe estar ligado a una sucursal y a una categoría. Por producto se
pueden definir diferentes variantes (tamaños) con precios individuales.

\subsubsection{Alta de Usuarios con Roles}

Roles participantes: Admin, Gerente. El usuario creador ingresa el correo del
nuevo usuario; el sistema genera una contraseña temporal. Al primer login, el
sistema solicita al nuevo usuario crear una contraseña propia. Las reglas de
negocio aplicables se detallan en la Tabla~\ref{tab:alta_usuarios}.

\subsubsection{Creación de Mesa Física}

Roles participantes: Mesero, Gerente. Desde la pantalla de visualización de
mesas, el usuario puede crear una nueva mesa física y asignarle un número según
la organización interna del restaurante.

\subsubsection{Inicialización de Mesa Virtual}

Roles participantes: Mesero, Gerente. Desde la pantalla de Mesas, el usuario
selecciona una o más mesas físicas en estado LIBRE. El sistema crea la mesa
virtual, la vincula a las mesas físicas y genera el QR con JWT. La mesa virtual
queda en estado PENDIENTE\_ACTIVACION. RB-M01: solo se pueden asociar mesas
físicas en estado LIBRE.

\subsubsection{Acceso de Comensales a la Mesa Virtual}

Roles participantes: Anfitrión, Comensal. Existen dos sub-flujos:

\noindent\textbf{Flujo Anfitrión — Activar Mesa:} El usuario selecciona
``Activar mi mesa''; el sistema activa la cámara para escanear el QR y valida
el JWT (válido, no expirado, no usado previamente). El anfitrión crea y confirma
una contraseña; la mesa pasa a estado POR\_LIQUIDAR. El sistema redirige al Home
de la mesa mostrando el ID y contraseña. Restricción: si dos usuarios intentan
activar la misma mesa simultáneamente, el primero en completar el flujo queda
como Anfitrión.

\noindent\textbf{Flujo Comensal — Unirse a Mesa:} El usuario selecciona
``Unirme a una mesa'', ingresa el identificador y contraseña. El sistema valida
las credenciales y redirige al Home de la mesa.

\subsubsection{Creación y Solicitud de Órdenes}

Roles participantes: Anfitrión, Comensal. El comensal navega el menú, selecciona
productos y los agrega a su canasta. Al finalizar su selección, elige entre:
(1)~\textbf{Ordenar individual} — los productos se cargan a su cuenta individual;
o (2)~\textbf{Agregar a mesa compartida} — los productos se suman a la cuenta
compartida de la mesa. Ambas acciones requieren un mensaje de confirmación explícito.
Se pueden generar $n$ órdenes individuales o $n$ adiciones a la mesa compartida
mientras la mesa esté en estado POR\_LIQUIDAR.

\subsubsection{Flujo de Liquidación}

Roles participantes: Anfitrión, Comensal, Mesero. El comensal accede a
``Proceder al pago'' y visualiza el desglose de cuentas. Selecciona qué cuentas
individuales desea liquidar y/o el monto a aportar a la mesa compartida. Tras el
pago exitoso, las cuentas seleccionadas se marcan como LIQUIDADAS. Cuando todos
los comensales y la mesa compartida estén en estado LIQUIDADA, la mesa virtual
cambia automáticamente a LIQUIDADA. El Mesero o Gerente debe liberar manualmente
la mesa (LIQUIDADA $\rightarrow$ LIBRE) para que las mesas físicas vuelvan a estar
disponibles.

% -----------------------------------------------------------
\subsection{Máquinas de Estado}
% -----------------------------------------------------------

Las siguientes máquinas de estado definen formalmente los ciclos de vida de las
entidades principales del sistema. Las transiciones son controladas por la capa
de API y son atómicas a nivel de base de datos.

\subsubsection{Estados de Mesa Virtual}

La Tabla~\ref{tab:estados_mesa} describe los estados posibles de una mesa virtual
y su semántica operacional.

\begin{table}[ht]
\caption{Estados de Mesa Virtual}
\label{tab:estados_mesa}
\centering
\small
\setlength{\tabcolsep}{5pt}
\begin{tabular}{>{\bfseries\raggedright\arraybackslash}p{4cm}
                >{\raggedright\arraybackslash}p{9cm}
                >{\centering\arraybackslash}p{3cm}}
\toprule
\thead{Estado} & \thead{Descripción} & \thead{Semántica} \\
\midrule
\rowcolor{orange}
PENDIENTE\_ACTIVACION &
  Mesa inicializada por Mesero/Gerente. QR generado. Aún no activada por el Anfitrión. &
  En espera \\
\rowcolor{green}
POR\_LIQUIDAR &
  Anfitrión escaneó el QR y creó contraseña. Comensales pueden ordenar; la mesa acumula deuda. &
  Activa \\
\rowcolor{lightblue}
LIQUIDADA &
  Todos los comensales y la cuenta compartida están en estado LIQUIDADA. No se aceptan más órdenes. &
  Completada \\
\rowcolor{altrow}
LIBRE &
  Liberada manualmente por Mesero o Gerente. Las mesas físicas asociadas regresan a estado LIBRE. &
  Disponible \\
\bottomrule
\end{tabular}
\end{table}

Las transiciones entre estados de la mesa virtual se detallan en la
Tabla~\ref{tab:trans_mesa}. Cabe destacar que cuando la mesa virtual se encuentra
en estado LIBRE, no se crea automáticamente una nueva mesa virtual; la creación
de una nueva mesa virtual solo ocurre cuando el Mesero o Gerente la inicializa
manualmente.

\begin{table}[ht]
\caption{Transiciones de Estado — Mesa Virtual}
\label{tab:trans_mesa}
\centering
\small
\setlength{\tabcolsep}{3pt}
\begin{tabular}{>{\raggedright\arraybackslash}p{3cm}
                >{\raggedright\arraybackslash}p{4cm}
                >{\raggedright\arraybackslash}p{3cm}
                >{\raggedright\arraybackslash}p{2.5cm}
                >{\raggedright\arraybackslash}p{3.5cm}}
\toprule
\thead{Estado Origen} & \thead{Evento / Condición} & \thead{Estado Destino} &
\thead{Actor} & \thead{Regla de Negocio} \\
\midrule
(Inexistente) &
  Mesero/Gerente inicializa mesa con mesas físicas libres &
  PENDIENTE\_ACTIVACION & Mesero / Gerente &
  RB-M01: mesas físicas en estado LIBRE \\
\rowcolor{altrow}
PENDIENTE\_ACTIVACION &
  Anfitrión escanea QR válido y crea contraseña &
  POR\_LIQUIDAR & Anfitrión &
  JWT válido, no expirado. Primera activación (usado = false) \\
POR\_LIQUIDAR &
  Todos los comensales y mesa compartida en LIQUIDADA &
  LIQUIDADA & Sistema (automático) &
  RB-L01: todos los pagos completados. Trigger en transacciones\_pago. \\
\rowcolor{altrow}
LIQUIDADA &
  Mesero o Gerente libera la mesa manualmente &
  LIBRE & Mesero / Gerente &
  RB-M03: acción manual obligatoria post-liquidación \\
LIBRE &
  Registro histórico — el registro no se reutiliza &
  (Archivado) & Sistema &
  Se crea nueva mesa virtual solo al inicializarla manualmente \\
\bottomrule
\end{tabular}
\end{table}

\subsubsection{Estados de Orden}

La Tabla~\ref{tab:estados_orden} describe los estados posibles de una orden
dentro del sistema.

\begin{table}[ht]
\caption{Estados de Orden}
\label{tab:estados_orden}
\centering
\small
\setlength{\tabcolsep}{5pt}
\begin{tabular}{>{\bfseries\raggedright\arraybackslash}p{3.5cm}
                >{\raggedright\arraybackslash}p{9.5cm}
                >{\centering\arraybackslash}p{3cm}}
\toprule
\thead{Estado} & \thead{Descripción} & \thead{Semántica} \\
\midrule
\rowcolor{orange}
PENDIENTE &
  Orden recibida por el sistema. Pendiente de confirmación por el Chef. & Esperando \\
\rowcolor{lightblue}
EN\_PREPARACION &
  Chef confirmó la orden. Se está preparando en cocina. & En proceso \\
\rowcolor{purple}
EN\_REPARTO &
  Chef marcó la preparación como finalizada. En espera de entrega. & Lista \\
\rowcolor{green}
ENTREGADA &
  El Comensal confirma haber recibido sus productos. & Completada \\
\rowcolor{red}
CANCELADA &
  Orden cancelada. Solo posible en estado PENDIENTE. & Cancelada \\
\bottomrule
\end{tabular}
\end{table}

\noindent\textit{Nota: La entrega de la orden al comensal debe ser marcada por
el Comensal (no por el Mesero), ya que es quien verifica haber recibido sus
productos.}

\medskip

La Tabla~\ref{tab:trans_orden} describe las transiciones de estado de una orden.

\begin{table}[ht]
\caption{Transiciones de Estado — Orden}
\label{tab:trans_orden}
\centering
\small
\setlength{\tabcolsep}{3pt}
\begin{tabular}{>{\raggedright\arraybackslash}p{3cm}
                >{\raggedright\arraybackslash}p{4cm}
                >{\raggedright\arraybackslash}p{3cm}
                >{\raggedright\arraybackslash}p{2.5cm}
                >{\raggedright\arraybackslash}p{3.5cm}}
\toprule
\thead{Estado Origen} & \thead{Evento / Condición} & \thead{Estado Destino} &
\thead{Actor} & \thead{Regla de Negocio} \\
\midrule
(Canasta) &
  Comensal confirma la orden & PENDIENTE & Comensal / Anfitrión &
  RB-O04: confirmación explícita. Mesa en estado POR\_LIQUIDAR. \\
\rowcolor{altrow}
PENDIENTE &
  Chef confirma inicio de preparación & EN\_PREPARACION & Chef &
  Tablero de Chef en tiempo real (Supabase Realtime). \\
EN\_PREPARACION &
  Chef marca la orden como lista & EN\_REPARTO & Chef &
  Notifica al Mesero. Transición visible en tiempo real. \\
\rowcolor{altrow}
EN\_REPARTO &
  Comensal confirma recepción de productos & ENTREGADA & Comensal &
  Transición irreversible. El Comensal valida la entrega. \\
PENDIENTE &
  Comensal/Mesero/Gerente cancela & CANCELADA & Mesero / Gerente &
  Solo cancelable en PENDIENTE. Los ítems se revierten del subtotal. \\
\bottomrule
\end{tabular}
\end{table}

\subsubsection{Estados de Cuenta Individual / Compartida}

La Tabla~\ref{tab:estados_cuenta} describe los estados posibles de una cuenta,
ya sea individual por comensal o la cuenta compartida de la mesa.

\begin{table}[ht]
\caption{Estados de Cuenta Individual / Compartida}
\label{tab:estados_cuenta}
\centering
\small
\setlength{\tabcolsep}{5pt}
\begin{tabular}{>{\bfseries\raggedright\arraybackslash}p{3.5cm}
                >{\raggedright\arraybackslash}p{9.5cm}
                >{\centering\arraybackslash}p{3cm}}
\toprule
\thead{Estado} & \thead{Descripción} & \thead{Semántica} \\
\midrule
\rowcolor{orange}
POR\_LIQUIDAR &
  Cuenta activa con saldo pendiente de liquidación. & Con deuda \\
\rowcolor{green}
LIQUIDADA &
  Monto total cubierto completamente por una o más transacciones. & Saldada \\
\bottomrule
\end{tabular}
\end{table}

La Tabla~\ref{tab:trans_cuenta} describe las transiciones de estado de una cuenta.

\begin{table}[ht]
\caption{Transiciones de Estado — Cuenta Individual / Compartida}
\label{tab:trans_cuenta}
\centering
\small
\setlength{\tabcolsep}{3pt}
\begin{tabular}{>{\raggedright\arraybackslash}p{3cm}
                >{\raggedright\arraybackslash}p{4cm}
                >{\raggedright\arraybackslash}p{3cm}
                >{\raggedright\arraybackslash}p{3cm}
                >{\raggedright\arraybackslash}p{3cm}}
\toprule
\thead{Estado Origen} & \thead{Evento / Condición} & \thead{Estado Destino} &
\thead{Actor} & \thead{Regla de Negocio} \\
\midrule
(Creación) &
  Primera orden cargada a la cuenta & POR\_LIQUIDAR & Sistema (trigger) &
  La cuenta se crea al generarse la primera orden. Monto calculado automáticamente. \\
\rowcolor{altrow}
POR\_LIQUIDAR &
  Pago procesa monto suficiente para cubrir el total & LIQUIDADA &
  Sistema (gateway webhook) &
  RB-L03/L04: control de concurrencia con SELECT FOR UPDATE. Transición atómica. \\
\bottomrule
\end{tabular}
\end{table}

% -----------------------------------------------------------
\subsection{Diseño de API}
% -----------------------------------------------------------

La API sigue convenciones REST con respuestas en JSON. Todos los endpoints
requieren el header \texttt{Authorization: Bearer \{jwt\}} con el token de
Supabase Auth, excepto los endpoints de sesión de comensal que utilizan el ID
de sesión.

\noindent\textit{Nota técnica: Este proyecto utiliza Server Actions de Next.js
(\texttt{"use server"}) como mecanismo de comunicación cliente-servidor. Las
funciones listadas en la Sección~\ref{sec:server_actions} son las implementaciones
actuales. Los endpoints REST que se describen a continuación definen el contrato
objetivo de la API pública.}

\subsubsection{Módulo: Mesas}

La Tabla~\ref{tab:api_mesas} detalla los endpoints del módulo de gestión de mesas.

\begin{table}[ht]
\caption{Endpoints REST — Módulo Mesas}
\label{tab:api_mesas}
\centering
\small
\setlength{\tabcolsep}{3pt}
\begin{tabular}{>{\centering\arraybackslash}p{1.5cm}
                >{\raggedright\arraybackslash}p{5cm}
                >{\raggedright\arraybackslash}p{3.5cm}
                >{\raggedright\arraybackslash}p{6cm}}
\toprule
\thead{Método} & \thead{Endpoint} & \thead{Auth requerida} & \thead{Descripción} \\
\midrule
GET &
  \texttt{/api/sucursales/:id/mesas} & Mesero, Gerente, Admin &
  Lista mesas físicas de la sucursal con estatus y mesa virtual asociada. \\
\rowcolor{altrow}
POST &
  \texttt{/api/mesas-fisicas} & Gerente, Admin &
  Crea una nueva mesa física en la sucursal. \\
POST &
  \texttt{/api/mesas-virtuales} & Mesero, Gerente, Admin &
  Inicializa una mesa virtual. Body: \texttt{\{mesa\_fisica\_ids[]\}}. Devuelve el token JWT para el QR. \\
\rowcolor{altrow}
POST &
  \texttt{/api/mesas-virtuales/:id/activar} & Pública (con QR token) &
  Activa la mesa y establece la contraseña. Body: \texttt{\{token\_qr, password\}}. Devuelve \texttt{sesion\_id}. \\
POST &
  \texttt{/api/mesas-virtuales/:id/unirse} & Pública &
  Comensal se une con ID y contraseña. Body: \texttt{\{codigo\_publico, password\}}. Devuelve \texttt{sesion\_id}. \\
\rowcolor{altrow}
PATCH &
  \texttt{/api/mesas-virtuales/:id/liberar} & Mesero, Gerente, Admin &
  Libera la mesa (LIQUIDADA $\rightarrow$ LIBRE). Solo si estatus = LIQUIDADA. \\
GET &
  \texttt{/api/mesas-virtuales/:id/estatus} & Sesión comensal o Empleado &
  Retorna estatus actual, comensales activos y montos de la mesa. \\
\bottomrule
\end{tabular}
\end{table}

\subsubsection{Módulo: Órdenes}

La Tabla~\ref{tab:api_ordenes} detalla los endpoints del módulo de órdenes.

\begin{table}[ht]
\caption{Endpoints REST — Módulo Órdenes}
\label{tab:api_ordenes}
\centering
\small
\setlength{\tabcolsep}{3pt}
\begin{tabular}{>{\centering\arraybackslash}p{1.5cm}
                >{\raggedright\arraybackslash}p{5cm}
                >{\raggedright\arraybackslash}p{3.5cm}
                >{\raggedright\arraybackslash}p{6cm}}
\toprule
\thead{Método} & \thead{Endpoint} & \thead{Auth requerida} & \thead{Descripción} \\
\midrule
GET &
  \texttt{/api/sucursales/:id/menu} & Sesión comensal &
  Menú completo de la sucursal con categorías, productos y variantes activas. \\
\rowcolor{altrow}
POST &
  \texttt{/api/ordenes} & Sesión comensal &
  Crea una orden. Body: \texttt{\{sesion\_id, tipo, items:[...]\}}. Valida mesa activa. \\
GET &
  \texttt{/api/mesas-virtuales/:id/ordenes} & Mesero, Chef, Gerente &
  Lista todas las órdenes de la mesa con sus ítems y estatus. \\
\rowcolor{altrow}
PATCH &
  \texttt{/api/ordenes/:id/estatus} & Chef, Mesero, Gerente &
  Actualiza estatus de la orden. Body: \texttt{\{estatus\}}. Solo transiciones válidas. \\
DELETE &
  \texttt{/api/ordenes/:id} & Mesero, Gerente &
  Cancela una orden en estado PENDIENTE. Revierte montos de las cuentas. \\
\bottomrule
\end{tabular}
\end{table}

\subsubsection{Módulo: Pagos}

La Tabla~\ref{tab:api_pagos} detalla los endpoints del módulo de pagos.

\begin{table}[ht]
\caption{Endpoints REST — Módulo Pagos}
\label{tab:api_pagos}
\centering
\small
\setlength{\tabcolsep}{3pt}
\begin{tabular}{>{\centering\arraybackslash}p{1.5cm}
                >{\raggedright\arraybackslash}p{5cm}
                >{\raggedright\arraybackslash}p{3.5cm}
                >{\raggedright\arraybackslash}p{6cm}}
\toprule
\thead{Método} & \thead{Endpoint} & \thead{Auth requerida} & \thead{Descripción} \\
\midrule
GET &
  \texttt{/api/mesas-virtuales/:id/liquidacion} & Sesión comensal o Empleado &
  Resumen completo: cuentas individuales, cuenta compartida, estatus por comensal. \\
\rowcolor{altrow}
POST &
  \texttt{/api/pagos/intento} & Sesión comensal &
  Crea intento de pago. Body: \texttt{\{sesion\_id, cuentas\_a\_liquidar[], monto\_compartida, idempotency\_key\}}. Retorna \texttt{client\_secret}. \\
POST &
  \texttt{/api/pagos/confirmar} & Webhook del gateway &
  Confirma el pago completado. Actualiza cuentas y transiciones de estado de mesa. \\
\bottomrule
\end{tabular}
\end{table}

\subsubsection{Módulo: Dashboard Analytics}

La Tabla~\ref{tab:api_analytics} detalla los endpoints del módulo de analytics.

\begin{table}[ht]
\caption{Endpoints REST — Módulo Dashboard Analytics}
\label{tab:api_analytics}
\centering
\small
\setlength{\tabcolsep}{3pt}
\begin{tabular}{>{\centering\arraybackslash}p{1.5cm}
                >{\raggedright\arraybackslash}p{5cm}
                >{\raggedright\arraybackslash}p{3.5cm}
                >{\raggedright\arraybackslash}p{6cm}}
\toprule
\thead{Método} & \thead{Endpoint} & \thead{Auth requerida} & \thead{Descripción} \\
\midrule
GET &
  \texttt{/api/analytics/sucursal/:id} & Gerente Suc., Gerente Zona, Admin &
  Métricas de la sucursal: mesas activas, ingresos, top productos, tiempo promedio. \\
\rowcolor{altrow}
GET &
  \texttt{/api/analytics/zona/:id} & Gerente Zona, Admin &
  Métricas agregadas de todas las sucursales de la zona. \\
GET &
  \texttt{/api/analytics/cadena/:id} & Admin &
  Métricas consolidadas de toda la cadena. \\
\bottomrule
\end{tabular}
\end{table}

\subsubsection{Server Actions Implementadas (Next.js)}
\label{sec:server_actions}

Las siguientes son las funciones de servidor actualmente implementadas mediante
Server Actions de Next.js. No son URLs públicas; son funciones invocadas
directamente por el runtime. La Tabla~\ref{tab:actions} presenta el inventario
completo por archivo.

\setstretch{1.2}
\setlength{\tabcolsep}{4pt}
\begin{longtable}{>{\ttfamily\raggedright\arraybackslash}p{4cm}
                >{\ttfamily\raggedright\arraybackslash}p{4cm}
                >{\raggedright\arraybackslash}p{8cm}}
\caption{Inventario de Server Actions por Archivo} \label{tab:actions} \\
\toprule
\thead{Archivo} & \thead{Función} & \thead{Descripción} \\
\midrule
\endfirsthead
\caption[]{TABLA \thetable{} (Continuación)} \\
\toprule
\thead{Archivo} & \thead{Función} & \thead{Descripción} \\
\midrule
\endhead
\bottomrule
\endfoot
chain.ts & getChainDashboard & Datos del dashboard de cadena (zonas, restaurantes, estadísticas). \\
\rowcolor{altrow}
chain.ts & getZoneDashboard & Datos del dashboard de la primera zona encontrada. \\
restaurant.ts & getDefaultRestaurant & Obtiene o crea el restaurante por defecto (demo). Envuelto en \texttt{React.cache()}. \\
\rowcolor{altrow}
menu.ts & getMenuData & Categorías e ítems del menú; crea categorías por defecto si no hay ninguna. \\
menu.ts & createCategory & Crea categoría (upsert por nombre + restaurante). \\
\rowcolor{altrow}
menu.ts & toggleItemSoldOut & Alterna estado agotado de un ítem. \\
menu.ts & deleteMenuItem & Elimina un ítem del menú. \\
\rowcolor{altrow}
menu.ts & updateMenuItem & Actualiza un ítem existente. \\
menu.ts & createMenuItem & Crea un ítem nuevo. \\
\rowcolor{altrow}
staff.ts & getStaffData & Lista personal; si está vacío, crea registros por defecto. \\
staff.ts & deleteStaffMember & Elimina un miembro del staff. \\
\rowcolor{altrow}
staff.ts & createStaffMember & Crea un miembro del staff. \\
staff.ts & toggleStaffStatus & Activa o desactiva un miembro. \\
\rowcolor{altrow}
orders.ts & getLiveOrders & Órdenes del día para la vista cocina (KDS). \\
orders.ts & advanceOrderStatus & Avanza el estado de la orden en el flujo estándar. \\
\rowcolor{altrow}
orders.ts & undoOrderStatus & Revierte el estado de la orden un paso. \\
orders.ts & moveOrderToStatus & Coloca la orden en un estado concreto. \\
\rowcolor{altrow}
tables.ts & getTables & Lista mesas del restaurante por defecto. \\
tables.ts & createTable & Crea mesa con QR y posición en rejilla. \\
\rowcolor{altrow}
tables.ts & updateTableStatus & Actualiza estado de mesa (vista dashboard/mesas). \\
tables.ts & deleteTable & Elimina una mesa. \\
\rowcolor{altrow}
tables.ts & updateTablePositions & Persiste posiciones (y forma opcional) en el plano. \\
waiter.ts & updateTableStatus & Actualiza estado de mesa; en DISPONIBLE rota QR y limpia código de unión. \\
\rowcolor{altrow}
waiter.ts & closeTable & Cierra sesiones activas y marca mesa como SUCIA. \\
waiter.ts & getWaiterTablesSummary & Resumen de mesas para la vista mesero. \\
\rowcolor{altrow}
waiter.ts & getTableDetail & Detalle de mesa: sesión, órdenes, total. \\
waiter.ts & waiterCreateOrder & Crea orden desde mesero (mesa ocupada y sesión activa). \\
\rowcolor{altrow}
waiter.ts & getMenuForOrdering & Menú para tomar pedidos (excluye agotados). \\
reports.ts & getDashboardReports & Estadísticas, top ítems y series para Hoy / Semana / Mes. \\
\rowcolor{altrow}
comensal.ts & submitComensalOrder & Envía pedido desde el flujo comensal (QR). \\
comensal.ts & getTableBill & Obtiene desglose de cuenta por comensal. \\
\rowcolor{altrow}
comensal.ts & guestJoinTable & Une comensal a la mesa; establece cookies. \\
comensal.ts & payGuestShare & Registra pago parcial de un comensal. \\
\rowcolor{altrow}
comensal.ts & requestBillAndPay & Pago, cierre de sesión y mesa en estado sucio. \\
comensal.ts & transferHost & Transfiere rol de anfitrión entre comensales. \\
\rowcolor{altrow}
comensal.ts & requestBill & Anfitrión solicita la cuenta (estado CERRANDO + broadcast). \\
comensal.ts & getGuestTableState & Estado para UI invitado (anfitrión, cuenta pedida, lista, código). \\
\rowcolor{altrow}
comensal.ts & getGuestOrders & Órdenes de la mesa visibles para el comensal. \\
\end{longtable}
\setstretch{1.5}

\subsubsection{Servicios Externos}

La Tabla~\ref{tab:servicios} describe los servicios externos utilizados por
la plataforma.

\begin{table}[ht]
\caption{Servicios Externos}
\label{tab:servicios}
\centering
\small
\setlength{\tabcolsep}{5pt}
\begin{tabular}{>{\bfseries\raggedright\arraybackslash}p{4cm}
                >{\raggedright\arraybackslash}p{12cm}}
\toprule
\thead{Servicio} & \thead{Uso} \\
\midrule
PostgreSQL &
  Base de datos principal vía Prisma (\texttt{DATABASE\_URL}). \\
\rowcolor{altrow}
Supabase &
  Cliente SSR y canales en tiempo real
  (\texttt{NEXT\_PUBLIC\_SUPABASE\_URL}, \texttt{NEXT\_PUBLIC\_SUPABASE\_ANON\_KEY}).
  Broadcast de órdenes, estado de mesa y liquidaciones. \\
\bottomrule
\end{tabular}
\end{table}

% -----------------------------------------------------------
\subsection{Diseño de Interfaces de Usuario (UI/UX)}
% -----------------------------------------------------------

El diseño de interfaces sigue tres principios rectores: \textit{Progressive
Disclosure} (revelar información según el contexto del usuario),
\textit{Mobile-First} (diseñado primero para pantallas de 375\,px y adaptado
para desktop) y accesibilidad WCAG 2.1 nivel AA (contraste mínimo 4.5:1 para
texto normal). La Tabla~\ref{tab:principios_ui} resume los principios de diseño
y su aplicación en la plataforma.

\begin{table}[ht]
\caption{Principios de Diseño de la Interfaz}
\label{tab:principios_ui}
\centering
\small
\setlength{\tabcolsep}{3pt}
\begin{tabular}{>{\bfseries\raggedright\arraybackslash}p{3.5cm}
                >{\raggedright\arraybackslash}p{6.0cm}
                >{\raggedright\arraybackslash}p{6.5cm}}
\toprule
\thead{Principio} & \thead{Descripción} & \thead{Aplicación en la plataforma} \\
\midrule
Mobile-First &
  Diseño base para pantallas de 375\,px, escalado vía breakpoints de Tailwind CSS. &
  Comensales usan el sistema exclusivamente desde smartphone. Chef y mesero desde tablet. Dashboard en desktop. \\
\rowcolor{altrow}
Progressive Disclosure &
  Solo se muestra la información relevante para el contexto actual del usuario. &
  El comensal ve el menú primero, sin información de otros comensales ni de administración. El chef solo ve el tablero. \\
Optimistic UI &
  La interfaz responde de inmediato a las acciones antes de recibir confirmación del servidor. &
  Agregar un producto a la canasta es inmediato. Si el servidor rechaza, la interfaz revierte con un mensaje de error claro. \\
\rowcolor{altrow}
Feedback Visual Claro &
  Cada acción tiene un estado de carga, éxito y error claramente distinguibles. &
  Botón ``Ordenar'' muestra spinner durante la petición. Estado de la orden visible en tiempo real. \\
Accesibilidad WCAG 2.1 AA &
  Contraste mínimo 4.5:1 para texto normal, 3:1 para texto grande. Navegación por teclado. ARIA labels. &
  Crítico para el tablero del Chef en condiciones de baja iluminación en cocina. \\
\bottomrule
\end{tabular}
\end{table}

\subsubsection{Rutas de Página}

La Tabla~\ref{tab:rutas} presenta el inventario de rutas de página del sistema
junto con su descripción.

\setstretch{1.2}
\setlength{\tabcolsep}{5pt}
\begin{longtable}{>{\ttfamily\raggedright\arraybackslash}p{6cm}
                >{\raggedright\arraybackslash}p{10cm}}
\caption{Rutas de Página del Sistema} \label{tab:rutas} \\
\toprule
\thead{Ruta} & \thead{Descripción} \\
\midrule
\endfirsthead
\caption[]{TABLA \thetable{} (Continuación)} \\
\toprule
\thead{Ruta} & \thead{Descripción} \\
\midrule
\endhead
\bottomrule
\endfoot
/ & Landing / Bienvenida \\
\rowcolor{altrow}
/admin & Administración \\
/barra & Vista barra \\
\rowcolor{altrow}
/cadena & Dashboard cadena \\
/cocina & KDS / Cocina \\
\rowcolor{altrow}
/dashboard & Dashboard principal \\
/dashboard/menu & Gestión de menú \\
\rowcolor{altrow}
/dashboard/mesas & Mesas (plano) \\
/dashboard/reportes & Reportes \\
\rowcolor{altrow}
/dashboard/settings & Ajustes \\
/dashboard/staff & Personal \\
\rowcolor{altrow}
/mesa/{[}codigo{]} & Entrada por QR de mesa \\
/mesa/{[}codigo{]}/menu & Menú comensal \\
\rowcolor{altrow}
/mesa/{[}codigo{]}/cuenta & Cuenta comensal \\
/mesero & Vista mesero \\
\rowcolor{altrow}
/roles & Roles \\
/zona & Vista zona \\
\end{longtable}
\setstretch{1.5}

\subsubsection{Mapa de Navegación por Rol}

A continuación se detalla la jerarquía de pantallas accesibles para cada rol.

\noindent\textbf{Mesero.} La Tabla~\ref{tab:nav_mesero} describe las pantallas
del flujo del Mesero.

\begin{table}[ht]
\caption{Navegación — Rol Mesero}
\label{tab:nav_mesero}
\centering
\small
\setlength{\tabcolsep}{3pt}
\begin{tabular}{>{\ttfamily\raggedright\arraybackslash}p{4cm}
                >{\raggedright\arraybackslash}p{6cm}
                >{\raggedright\arraybackslash}p{6cm}}
\toprule
\thead{Ruta} & \thead{Descripción} & \thead{Flujo de navegación} \\
\midrule
/login &
  Formulario de correo y contraseña. Redirección a cambio de contraseña en primer login. &
  Punto de entrada para empleados. \\
\rowcolor{altrow}
/mesas &
  Vista principal: grid de mesas físicas con estatus en tiempo real y color semántico. &
  Home del mesero tras login. \\
/mesas/nueva &
  Formulario para dar de alta una nueva mesa física. &
  Accesible desde botón '+' en /mesas. \\
\rowcolor{altrow}
/mesas/:id/inicializar &
  Selección de mesas físicas libres para crear una nueva mesa virtual. Genera el QR. &
  Accesible desde una mesa en estado LIBRE. \\
/mesas/:id/detalle &
  Detalle de la mesa virtual: estatus, color semántico, código, comensales,
  monto en deuda, monto liquidado. &
  Accesible desde una mesa activa. \\
\bottomrule
\end{tabular}
\end{table}

\noindent\textbf{Comensales (Anfitrión / Comensal).} La Tabla~\ref{tab:nav_comensal}
describe las pantallas del flujo de comensales.

\begin{table}[ht]
\caption{Navegación — Roles Comensal y Anfitrión}
\label{tab:nav_comensal}
\centering
\small
\setlength{\tabcolsep}{3pt}
\begin{tabular}{>{\ttfamily\raggedright\arraybackslash}p{4cm}
                >{\raggedright\arraybackslash}p{6cm}
                >{\raggedright\arraybackslash}p{6cm}}
\toprule
\thead{Ruta} & \thead{Descripción} & \thead{Flujo de navegación} \\
\midrule
/ &
  Pantalla de bienvenida: ``Activar mi mesa'' (Anfitrión) o ``Unirme a una mesa'' (Comensal). &
  Punto de entrada para comensales. \\
\rowcolor{altrow}
/activar &
  Cámara QR (html5-qrcode) y formulario de nueva contraseña con confirmación. &
  Solo al elegir ``Activar mi mesa''. No se puede activar una mesa ya activada. \\
/unirse &
  Formulario para ingresar ID de mesa y contraseña. &
  Solo al elegir ``Unirme a una mesa''. \\
\rowcolor{altrow}
/mesa/:id &
  Home de la mesa: ID y contraseña visibles, accesos rápidos, estatus. &
  Destino tras activación o unión exitosa. \\
/mesa/:id/menu &
  Catálogo de productos por categoría, imágenes, variantes y botón ``Agregar a canasta''. &
  Accesible desde el home de la mesa. \\
\rowcolor{altrow}
/mesa/:id/canasta &
  Lista de productos con cantidades y total. Botones: ``Ordenar individual'' y ``Agregar a mesa compartida''. &
  Accesible desde el ícono de canasta en el header. \\
/mesa/:id/pagar &
  Desglose de cuentas: individual por comensal + mesa compartida. Selección de qué liquidar. &
  Accesible desde botón ``Proceder al pago''. \\
\rowcolor{altrow}
/mesa/:id/checkout &
  Formulario de pago con total a cobrar. Integración con Stripe Elements / gateway. &
  Destino tras confirmar selección en /pagar. \\
/mesa/:id/recibo &
  Comprobante digital: ítems liquidados, monto y timestamp. &
  Destino tras pago exitoso. \\
\bottomrule
\end{tabular}
\end{table}

\noindent\textbf{Chef.} La Tabla~\ref{tab:nav_chef} describe las pantallas del
flujo del Chef.

\begin{table}[ht]
\caption{Navegación — Rol Chef}
\label{tab:nav_chef}
\centering
\small
\setlength{\tabcolsep}{3pt}
\begin{tabular}{>{\ttfamily\raggedright\arraybackslash}p{4cm}
                >{\raggedright\arraybackslash}p{6cm}
                >{\raggedright\arraybackslash}p{6cm}}
\toprule
\thead{Ruta} & \thead{Descripción} & \thead{Flujo de navegación} \\
\midrule
/login &
  Formulario de acceso. Redirección a /cocina tras autenticación. &
  Punto de entrada. \\
\rowcolor{altrow}
/cocina &
  Tablero Kanban: columnas PENDIENTE / EN\_PREPARACION / EN\_REPARTO.
  Actualización en tiempo real. &
  Home del chef. Pantalla principal de trabajo. \\
/cocina/:orden\_id &
  Detalle de orden: ítems, cantidades, notas del comensal. Botones de cambio de estatus. &
  Accesible al hacer clic en una orden del tablero. \\
\bottomrule
\end{tabular}
\end{table}

\noindent\textbf{Gerente y Admin.} La Tabla~\ref{tab:nav_gerente} describe las
pantallas del flujo de Gerentes y Administradores.

\begin{table}[ht]
\caption{Navegación — Roles Gerente y Admin}
\label{tab:nav_gerente}
\centering
\small
\setlength{\tabcolsep}{3pt}
\begin{tabular}{>{\ttfamily\raggedright\arraybackslash}p{4cm}
                >{\raggedright\arraybackslash}p{6cm}
                >{\raggedright\arraybackslash}p{6cm}}
\toprule
\thead{Ruta} & \thead{Descripción} & \thead{Roles con acceso} \\
\midrule
/dashboard &
  Analytics: métricas del período, gráficas de ingresos, top productos, mesas activas. &
  Gerente (filtrado por nivel), Admin (toda la cadena). \\
\rowcolor{altrow}
/zonas & CRUD de zonas. Solo visible con nivel Cadena. & Admin. \\
/sucursales & CRUD de sucursales. Filtrado por zonas asignadas. & Admin, Gerente Zona. \\
\rowcolor{altrow}
/menu & Gestión de categorías y productos. CRUD completo con activación/desactivación. &
  Admin, Gerente. \\
/usuarios &
  Alta, edición y desactivación de usuarios colaboradores (según RB-U04 a RB-U09). &
  Admin, Gerente. \\
\rowcolor{altrow}
/mesas & Vista de mesas igual que mesero, más botón de liberación manual. & Admin, Gerente. \\
\bottomrule
\end{tabular}
\end{table}

\subsubsection{Inventario de Pantallas}

La Tabla~\ref{tab:inventario} presenta el inventario completo de pantallas con
sus características técnicas principales.

\setstretch{1.2}
\setlength{\tabcolsep}{4pt}
\begin{longtable}{>{\raggedright\arraybackslash}p{3cm}
                >{\raggedright\arraybackslash}p{2cm}
                >{\ttfamily\raggedright\arraybackslash}p{3cm}
                >{\raggedright\arraybackslash}p{4.5cm}
                >{\centering\arraybackslash}p{1.5cm}}
\caption{Inventario de Pantallas y Características Técnicas} \label{tab:inventario} \\
\toprule
\thead{Pantalla} & \thead{Rol} & \thead{Ruta} &
\thead{Componentes principales} & \thead{Tiempo real} \\
\midrule
\endfirsthead
\caption[]{TABLA \thetable{} (Continuación)} \\
\toprule
\thead{Pantalla} & \thead{Rol} & \thead{Ruta} &
\thead{Componentes principales} & \thead{Tiempo real} \\
\midrule
\endhead
\bottomrule
\endfoot
Home comensales & Comensal / Anfitrión & / & Botones de acción, ilustración, logo. & No \\
\rowcolor{altrow}
Activar mesa & Anfitrión & /activar & Cámara QR (html5-qrcode), formulario contraseña. & No \\
Unirse a mesa & Comensal & /unirse & Formulario ID + contraseña, validación. & No \\
\rowcolor{altrow}
Home mesa & Comensal / Anfitrión & /mesa/:id & Badge ID+contraseña, accesos rápidos, estatus. & Sí \\
Menú restaurante & Comensal / Anfitrión & /mesa/:id/menu & Filtros por categoría, cards, variantes, canasta. & No \\
\rowcolor{altrow}
Canasta & Comensal / Anfitrión & /mesa/:id/canasta & Lista ítems, ajuste cantidad, total, botones. & No \\
Liquidación & Comensal / Anfitrión & /mesa/:id/pagar & Desglose por comensal, mesa compartida, checkboxes. & Sí \\
\rowcolor{altrow}
Checkout & Comensal / Anfitrión & /mesa/:id/checkout & Stripe Elements / widget gateway, total, pago. & No \\
Recibo digital & Comensal / Anfitrión & /mesa/:id/recibo & Resumen de pago, ítems, timestamp. & No \\
\rowcolor{altrow}
Login empleados & Todos empleados & /login & Formulario email/password. & No \\
Mesas físicas & Mesero / Gerente / Admin & /mesas & Grid con color semántico, panel detalle lateral. & Sí \\
\rowcolor{altrow}
Inicializar mesa virtual & Mesero / Gerente / Admin & /mesas/:id/inicializar & Selección mesas libres, preview QR. & No \\
Tablero Chef (Kanban) & Chef & /cocina & 3 columnas Kanban, cards con tiempo transcurrido. & Sí \\
\rowcolor{altrow}
Dashboard Analytics & Gerente / Admin & /dashboard & Cards KPIs, gráfica ingresos (Chart.js), top productos. & No \\
Gestión de menú & Gerente / Admin & /menu & Árbol categorías-productos, modal creación/edición. & No \\
\rowcolor{altrow}
Gestión de usuarios & Admin / Gerente & /usuarios & Tabla usuarios, formulario de alta, badge de rol. & No \\
\end{longtable}
\setstretch{1.5}

\subsubsection{Especificación de Componentes Clave}

\noindent\textbf{Tablero Kanban del Chef (/cocina).} La actualización en tiempo
real se implementa mediante Supabase Realtime suscrito al canal \texttt{ordenes},
filtrado por \texttt{sucursal\_id}; el tablero se actualiza sin recarga al cambiar
el estatus de cualquier orden. Al llegar una orden PENDIENTE, se reproduce un
sonido via Web Audio API y se muestra un badge en la columna correspondiente.
Cada card de orden muestra el número de mesa virtual, el tiempo transcurrido
desde la creación (contador en tiempo real), la lista de ítems con cantidades,
notas del comensal y el botón de cambio de estatus. El contador se actualiza
cada segundo (setInterval); el fondo cambia a amarillo al superar los 10\,min
y a rojo al superar los 20\,min. Para la accesibilidad en cocina, la fuente mínima
es de 16\,px, el contraste es $\geq 7$:1 en modo oscuro y el área táctil mínima
por botón es de 44$\times$44\,px.

\medskip

\noindent\textbf{Canasta y Flujo de Orden (/mesa/:id/canasta).} El estado de
la canasta se gestiona con Zustand (estado global del cliente), persiste durante
la sesión del navegador y se vacía al confirmar la orden. Al agregar un producto
con múltiples variantes, se presenta un modal de selección con precios en tiempo
real. La confirmación de orden requiere un dialog modal con resumen de ítems y
monto total; el usuario debe hacer clic explícito en ``Ordenar individual'' o
``Agregar a mesa compartida''. Antes de enviar la orden, el cliente verifica el
estatus de la mesa; si está en LIQUIDADA, el botón de confirmación se bloquea
y se muestra un mensaje de error.

\medskip

\noindent\textbf{Componente de Liquidación (/mesa/:id/pagar).} El desglose de
cuentas se presenta mediante un acordeón expandible por comensal, mostrando
órdenes individuales y subtotales. La cuenta compartida se muestra al final con
todos sus ítems y el comensal que los agregó. La selección de pago se realiza
mediante checkboxes para cuentas individuales e input numérico para la aportación
a la mesa compartida, con validación de rango $[0, \text{monto\_pendiente\_compartida}]$.
El botón ``Liquidar todo'' selecciona todo automáticamente. Supabase Realtime
actualiza el estatus de pago de cada comensal; las cuentas liquidadas se muestran
en verde con ícono de verificación. Para la prevención de doble pago, el botón
``Proceder al pago'' se deshabilita inmediatamente al hacer clic y la API genera
la \texttt{idempotency\_key} antes de llamar al gateway.

% ============================================================
%  CAPÍTULO 6 — IMPLEMENTACIÓN Y PRUEBAS  (en blanco)
% ============================================================
\newpage
